\section{A nonlinear Green formula from coherent moments}
\label{cy:section}

The product-law energy limit has a differential consequence. It
recovers the Calabi--Yau potential itself from random interpolation
sections, rather than only recovering its variational energy. The
representation below is the determinant formula proposed in
\cite[Conjecture 6.1]{prob_berman2017}. Its finite-level convexity
mechanism reduces a two-coordinate determinant to the projective line.
The underlying product-energy theorem addresses the determinant instance
of \cite[Problem 4.10]{prob_berman2017}, rather than arbitrary singular
Hamiltonians. A later conditional result obtained the same potential
formula from a uniform complex zero-free neighborhood of the partition
functions \cite[Theorems 3.2--3.3]{Berman2023zeta}. The argument here
uses the real-parameter energy estimate and finite-level convexity;
it makes no assertion about the complex zeros of those functions.
In complex dimension one the representation was already unconditional
by bosonization, as recalled in \cite[Remark 3.4]{Berman2023zeta}.
The argument below applies in every complex dimension under the stated
smooth positive hypotheses.

Let $X$ be a compact connected complex manifold of dimension $q\ge1$,
and let $(L,h)$ be a positive Hermitian line bundle with curvature
$\omega=dd^c\phi_0$, where $dd^c=i\partial\bar\partial/(2\pi)$.
Put $V=\int_X\omega^q$. For all sufficiently large $k$, let
$N=N_k=h^0(X,L^k)$ and choose an orthonormal basis for a fixed smooth
positive reference volume $\sigma$. Denote its evaluation determinant
by $\Delta_k$ and keep this basis fixed when varying a probability
measure $\mu$. As before,
\[
 E_k=-\frac1{kN}\log\|\Delta_k\|^2,
 \qquad \overline E_k(\mu)=\int E_k\,d\mu^{\otimes N}.
\]
All norms use the fixed metric $h$, in every determinant coordinate.

\begin{lemma}[Convexity of the product mean energy]
\label{cy:finite-convexity}
For every sufficiently large fixed $k$, the functional
$\overline E_k$ is convex on smooth strictly positive probability
volumes. If $\eta$ is a smooth signed measure of total mass zero and
$\mu_t=\mu+t\eta$ remains strictly positive, then
\begin{equation}
 \frac{d}{dt}\overline E_k(\mu_t)\bigg|_{t=0}
 =-\int_X U_{k,\mu}(x)\,d\eta(x),\qquad
 U_{k,\mu}(x)=\frac1k\int_{X^{N-1}}
        \log\|\Delta_k(x,\mathbf y)\|^2\,d\mu^{N-1}(\mathbf y).
 \label{cy:first-variation}
\end{equation}
\end{lemma}

\begin{proof}
The logarithm of a nonzero holomorphic section is locally integrable.
On the compact product manifold this makes the determinant logarithm
absolutely integrable against every product of smooth volume densities.
Thus expansion of $(\mu+t\eta)^N$ gives an ordinary polynomial in
$t$, and differentiation gives \eqref{cy:first-variation} and
\[
 \frac{d^2}{dt^2}\overline E_k(\mu_t)
 =-\frac{N-1}{k}\int_{X^{N-2}}
   \left[\int_{X^2}\log\|\Delta_k(x,y,\mathbf z)\|^2
                      \,d\eta(x)d\eta(y)\right]d\mu_t^{N-2}(\mathbf z).
\]
For almost every retained tuple $\mathbf z$, its evaluation rows have
rank $N-2$. Indeed, a nonzero full determinant gives one such tuple,
and rank failure is a proper analytic locus. The two-dimensional
annihilator of these rows is a pencil of sections $a,b$ of $L^k$.
Expansion of the determinant gives a nonzero constant times
$a(x)b(y)-b(x)a(y)$, with the retained norm factors constant in $x,y$.
Set
\[
 r(x)=\|a(x)\|^2+\|b(x)\|^2,\qquad
 P(x)=P_{[a(x):b(x)]}\in\mathbb P^1.
\]
The common zero set has volume zero. Off it,
\[
 \log\|\Delta_k(x,y,\mathbf z)\|^2
 =c_{\mathbf z}+\log r(x)+\log r(y)
                     +\log[1-\operatorname{tr}(P(x)P(y))].
\]
Since $\eta(X)=0$, the first three terms contribute zero to the signed
double integral. The remaining kernel is positive definite after
changing its sign:
\begin{equation}
 K(P,Q)=-\log[1-\operatorname{tr}(PQ)]
   =\sum_{j\ge1}\frac1j
            \operatorname{tr}(P^{\otimes j}Q^{\otimes j}).
 \label{cy:projective-kernel}
\end{equation}
Each finite sum has nonnegative quadratic form because
\[
 \iint\operatorname{tr}(P(x)^{\otimes j}P(y)^{\otimes j})
                         \,d\eta(x)d\eta(y)
 =\left\|\int P(x)^{\otimes j}\,d\eta(x)\right\|_{\rm HS}^2\ge0.
\]
To justify the infinite sum with a signed measure, first apply Fubini
to the absolute determinant logarithm. For almost every $\mathbf z$
its double integral against $|\eta|\otimes|\eta|$ is finite. The
functions $\log r$ are also integrable: they are bounded above, and a
nonzero component section bounds their negative parts. The displayed
kernel identity therefore makes $K$ integrable against
$|\eta|\otimes|\eta|$. Its increasing finite sums are bounded by $K$,
so dominated convergence proves nonnegativity of the full quadratic
form. More precisely, if $P_{\mathbf z}$ denotes the pencil map for
these retained rows, the computation proves the exact hierarchy
\begin{equation}
 \frac{d^2}{dt^2}\overline E_k(\mu_t)
 =\frac{N-1}{k}\int_{X^{N-2}}
   \sum_{j\ge1}\frac1j
   \left\|\int_X P_{\mathbf z}(x)^{\otimes j}\,d\eta(x)
          \right\|_{\rm HS}^2
                  \,d\mu_t^{N-2}(\mathbf z)\ge0.
 \label{cy:hessian-hierarchy}
\end{equation}
Tonelli justifies averaging the nonnegative terms over $\mathbf z$.
Thus coherent moment discrepancies form the finite-level Hessian.
No convergence of these Hessians is required for the result below.
\end{proof}

\begin{theorem}[Nonlinear Green representation]
\label{cy:green-formula}
Let $\mu$ be any smooth strictly positive probability volume on $X$.
There is a unique smooth potential $\phi_\mu$, normalized by
$\int_X\phi_\mu\,d\mu=0$, such that
\[
 \frac{(\omega+dd^c\phi_\mu)^q}{V}=\mu,
 \qquad \omega+dd^c\phi_\mu>0.
\]
The mean-zero determinant potentials
\begin{equation}
 \phi_{k,\mu}(x)=\frac1k\int_{X^{N_k-1}}
       \log\|\Delta_k(x,\mathbf y)\|^2\,d\mu^{N_k-1}(\mathbf y)
      -C_{k,\mu},
 \quad
 C_{k,\mu}=\frac1k\int_{X^{N_k}}\log\|\Delta_k\|^2\,d\mu^{N_k}
 \label{cy:potential-formula}
\end{equation}
converge to $\phi_\mu$ in $L^p(X,\sigma)$, equivalently in
$L^p(X,\mu)$, for every fixed $1\le p<\infty$. The formula is independent of the chosen basis after
the displayed centering.
\end{theorem}

\begin{proof}
Existence and uniqueness of the smooth solution are Yau's theorem
\cite{Yau1978}. The energy dual is attained at that solution:
\[
 E_\omega(\mu)=\mathcal A(\phi_\mu)-\int\phi_\mu\,d\mu.
\]
For a smooth signed variation $\eta$ of mass zero, the normalized
solutions for $\mu+t\eta$ vary smoothly for small $t$. This follows
from the elliptic implicit function theorem: the linearization of the
normalized Monge--Amp\`ere map is the Laplacian of the positive metric,
an isomorphism modulo constants onto zero-integral variations.
Differentiating the last energy identity and using
$d\mathcal A_\phi(v)=\int v\,\operatorname{MA}_\omega(\phi)$ gives
\begin{equation}
 \frac{d}{dt}E_\omega(\mu+t\eta)\bigg|_{t=0}
                 =-\int\phi_\mu\,d\eta.
 \label{cy:energy-derivative}
\end{equation}

The product mean-energy theorem
\ref{att:general-product-energy}, in its polarized form explained in
Corollary~\ref{att:polarized-extension}, gives
$\overline E_k(\mu+t\eta)\to E_\omega(\mu+t\eta)$ for every
fixed sufficiently small $t$. The functions of $t$ on the left are
convex by Lemma~\ref{cy:finite-convexity}. Hence the secant inequalities
\[
 \frac{\overline E_k(\mu)-\overline E_k(\mu-\varepsilon\eta)}\varepsilon
 \le \frac{d}{dt}\overline E_k(\mu+t\eta)\bigg|_{t=0}
 \le
 \frac{\overline E_k(\mu+\varepsilon\eta)-\overline E_k(\mu)}\varepsilon
\]
show, by first taking $k\to\infty$ and then $\varepsilon\downarrow0$,
that their derivatives converge to \eqref{cy:energy-derivative}.
Because constants pair to zero with $\eta$, this says
\[
 \int\phi_{k,\mu}\,d\eta\longrightarrow\int\phi_\mu\,d\eta
 \quad\text{for every smooth zero-mass signed }\eta.
\]

For completeness, these weak identities determine an $L^1$ limit.
The averaged logarithm in \eqref{cy:potential-formula} defines an
$L^1$ function, by absolute logarithmic integrability and Fubini.
In distributions it satisfies
$dd^c\phi_{k,\mu}\ge-\omega$, since each determinant slice is a
section of $L^k$. Its psh representative is therefore an
$\omega$-psh potential. Centering makes its $\mu$ mean zero.
The normalized-psh exponential estimate gives a uniform constant $B_\mu$
such that every $\omega$-psh $v$ with $\sup v=0$ satisfies
$\int(-v)\,d\mu\le B_\mu$. Consequently
\[
 0\le\sup_X\phi_{k,\mu}\le B_\mu.
\]
Compactness of normalized psh potentials makes $\phi_{k,\mu}$ relatively
compact in $L^1(\sigma)$. Every cluster point has mean zero against
$\mu$, and its integrals against all smooth zero-mass signed measures
agree with those of $\phi_\mu$. Their difference is constant as a
distribution, and normalization makes that constant zero. Every cluster
point is therefore $\phi_\mu$, proving $L^1$ convergence of the full
sequence. The same uniform exponential estimate, together with the
uniform upper bounds, makes $|\phi_{k,\mu}|^p$ uniformly integrable
for every finite $p$. Vitali therefore upgrades the convergence to
$L^p$. A fixed basis change multiplies the determinant by a nonzero constant;
its logarithm cancels exactly in \eqref{cy:potential-formula}.
\end{proof}

\begin{corollary}[Random interpolation divisors and sampling error]
\label{cy:interpolation-sampling}
Draw $\mathbf Y=(Y_2,\ldots,Y_N)$ independently from $\mu$.
Almost surely the section
$s_{\mathbf Y}(x)=\Delta_k(x,\mathbf Y)$ is nonzero and is the unique
section vanishing at these $N-1$ points, up to scale. Define its
mean-zero potential
\[
 u_{k,\mathbf Y}(x)=\frac1k\log\|s_{\mathbf Y}(x)\|^2
                 -\int_X\frac1k\log\|s_{\mathbf Y}\|^2\,d\mu.
\]
Then
\begin{equation}
 \mathbb E\frac1k[Z_{s_{\mathbf Y}}]
       =\omega+dd^c\phi_{k,\mu}
       \longrightarrow\omega+dd^c\phi_\mu
 \label{cy:divisor-limit}
\end{equation}
weakly as currents. Moreover, if $u^{(1)},\ldots,u^{(M)}$ are independent
copies of $u_{k,\mathbf Y}$, then, for every $a>0$ with
\[
 C_a(\mu):=\sup_{\substack{v\in\operatorname{PSH}(X,\omega)\\\sup v=0}}
                      \int e^{-av}\,d\mu<\infty,
\]
one has the uniform sampling estimate
\begin{equation}
 \mathbb E\left\|\frac1M\sum_{j=1}^M u^{(j)}-\phi_{k,\mu}
             \right\|_{L^2(\mu)}^2
       \le\frac{2C_a(\mu)}{a^2M}.
 \label{cy:sampling-bound}
\end{equation}
On the anticanonical projective plane, if $\mu=g\sigma$, one may take
$a=1/6$, obtaining the explicit bound $192\|g\|_\infty/M$.

Exact centering integrals can also be replaced by sampling. For each
section independently draw $J\ge1$ calibration points from $\mu$
and subtract their mean log squared norm in place of its $\mu$ integral.
Then the left side of \eqref{cy:sampling-bound} is bounded by
\begin{equation}
 \frac{2C_a(\mu)}{a^2M}\left(1+\frac1J\right).
 \label{cy:sample-centered-bound}
\end{equation}
\end{corollary}

\begin{proof}
Independence of $N-1$ evaluation rows fails only on a proper analytic
subset, so the interpolation assertion holds almost surely. The
Poincar\'e--Lelong formula in our convention is
$k^{-1}[Z_s]=\omega+dd^c(k^{-1}\log\|s\|^2)$.
Taking expectations and using \eqref{cy:potential-formula} proves the
identity of currents; $L^1$ convergence of potentials proves its limit.

For the quantitative statement subtract the supremum from the section
potential and write $v\le0$ for the result. Centering at its $\mu$ mean
gives $u=v-\mu(v)$, so
\[
 \|u\|_{L^2(\mu)}^2\le\int v^2\,d\mu
              \le\frac2{a^2}\int e^{-av}\,d\mu
              \le\frac{2C_a(\mu)}{a^2}.
\]
Here $x^2\le2a^{-2}e^{ax}$ for $x\ge0$ follows from the exponential
series. The random functions form a square-integrable random variable
in the separable Hilbert space $L^2(\mu)$, and their mean is
$\phi_{k,\mu}$. The elementary variance identity for independent
Hilbert-space samples gives \eqref{cy:sampling-bound}.
On $\mathbb P^2$, the scalar bound gives
$C_{1/6}(\mu)\le(8/3)\|g\|_\infty$, yielding the stated constant.

For the sample-centered version, write the calibration error for section
$j$ as $e_j=J^{-1}\sum_{\ell=1}^J u^{(j)}(Z_{j\ell})$.
Conditionally on the interpolation sections it has mean zero and
variance at most $2C_a(\mu)/(a^2J)$, and the $e_j$ are conditionally
independent. The final estimator is
$M^{-1}\sum_j u^{(j)}-M^{-1}\sum_j e_j$.
The second term is a random constant of squared $L^2(\mu)$ norm
$|M^{-1}\sum_j e_j|^2$. Its cross term with the first centered
estimation error has conditional mean zero. Adding variances gives
\eqref{cy:sample-centered-bound}.
\end{proof}

\paragraph{Computational access and scope.}
A direct sampling procedure takes a basis of $H^0(X,L^k)$ that can be
evaluated, the Hermitian norm, and independent sample access to $\mu$.
For each repetition, draw $N_k-1$ interpolation points, form their
evaluation matrix, and choose a nonzero null vector. Evaluate the
resulting section, subtract its mean log squared norm on $J$ fresh calibration
points, divide by $k$, and average over the $M$ repetitions. Changing
the normalization of any null vector cancels in this subtraction.
Thus no exact integration oracle is needed. The theorem specifies the
sampling error of this ideal arithmetic procedure, without an assertion
of efficient sample access or finite-precision conditioning.

The sampling bound controls the error at a fixed polynomial level.
Theorem~\ref{cy:green-formula} separately proves convergence of the
level-dependent bias, without an effective rate in $k$. Taking
$k\to\infty$ and $M\to\infty$, even with $J=1$, gives convergence
in expected squared $L^2(\mu)$ to the Calabi--Yau potential: the error
is the sampling variance plus $\|\phi_{k,\mu}-\phi_\mu\|_2^2$.
No uniform-norm or higher-derivative convergence is asserted.
If the canonical bundle is trivial and $\mu$ is the normalized volume
of a nowhere-vanishing holomorphic top form, the limiting K\"ahler
form is the Ricci-flat metric in $c_1(L)$. In particular the K3 surface
of Theorem~\ref{geo:k3-lattice}, equipped with any ample polarization,
is an algebraic test case for the construction. This application uses
its Calabi--Yau structure; it does not assert an effective polynomial-
level bias bound or a practical sampling complexity for that surface.
